\documentclass[12pt,a4paper]{article}

\usepackage[utf8]{inputenc}
\usepackage[T2A]{fontenc}
\usepackage[russian,english]{babel}
\usepackage{geometry}
\usepackage{longtable}
\usepackage{array}
\usepackage{enumitem}
\usepackage{hyperref}

\geometry{margin=2cm}

\hypersetup{
    colorlinks=true,
    linkcolor=blue,
    urlcolor=blue
}

\setlist[itemize]{topsep=4pt,itemsep=2pt}
\setlist[enumerate]{topsep=4pt,itemsep=2pt}

\begin{document}

\selectlanguage{russian}

\begin{center}
    {\LARGE \textbf{Sprint 1: Required Report Deliverables}}\\[6pt]
    {\large Обязательные материалы для вставки в LaTeX report}
\end{center}

\section{Общее правило сдачи материалов}

Каждый участник должен отправить только те материалы, которые нужны для блока \textbf{Deliverables (in Report)}.
Материалы можно сдавать в одном из двух форматов:

\begin{itemize}
    \item готовая LaTeX-секция, которую можно сразу вставить в report;
    \item ссылка на файл, pull request, issue, screenshot, diagram или API documentation.
\end{itemize}

Если участник не успевает дать какую-либо ссылку или артефакт, он должен отправить короткую причину отсутствия. Рома вставляет эту причину в report вместо ссылки.

\textbf{Формат причины отсутствия:}

\begin{verbatim}
Artifact is not available yet because ...
\end{verbatim}

Пример:

\begin{verbatim}
OpenAPI link is not available yet because backend endpoints are still
being finalized. The API contract will be added in the next sprint.
\end{verbatim}

\section{Минимальный обязательный набор для report}

В report обязательно должны попасть следующие элементы:

\begin{enumerate}
    \item Team name, members, assigned roles and communication channels.
    \item Chosen project, chosen track, problem statement and Industrial track justification.
    \item Project plan: roadmap, milestones, dependencies, risks and Kanban backlog link.
    \item User stories, acceptance criteria and MVP scope.
    \item Tech stack with justification.
    \item Wireframes/mockups and user flow diagrams.
    \item Description of implemented MVP features and functional user journey.
    \item Links to relevant PRs/Issues and API documentation, if available.
    \item ML experimentation setup and initial baseline result.
    \item Initial measurement against the baseline product.
    \item Internal demo notes and identified next steps.
\end{enumerate}

\section{Deliverables by Team Member}

\begin{longtable}{|p{0.16\textwidth}|p{0.74\textwidth}|}
\hline
\textbf{Участник} & \textbf{Что обязательно отправить для report} \\
\hline

\textbf{Данил} &
\begin{itemize}
    \item LaTeX-секция \texttt{Project Track Justification}: chosen project, Industrial track, explanation of how the project satisfies Industrial track criteria.
    \item LaTeX-секция \texttt{Project Plan}: roadmap на оставшиеся недели, milestones, dependencies and risks.
    \item LaTeX-секция \texttt{MVP Scope}: what is IN and what is OUT.
    \item Link на Kanban backlog.
    \item LaTeX-секция \texttt{Tech Stack Justification}.
    \item LaTeX-секция \texttt{Industrial Baseline Measurement}: как GitFlame работает сейчас и что наша система должна улучшить.
\end{itemize}
Если Kanban link или baseline measurement не готовы, отправить короткую причину отсутствия. \\
\hline

\textbf{Карим} &
\begin{itemize}
    \item LaTeX-секция или ссылка на файл \texttt{ML Experimentation Setup}: как тестируется open-source model или ML service.
    \item LaTeX-секция \texttt{Initial Baseline Result}: первый результат анализа или рекомендаций.
    \item Если есть prompt/schema files, дать ссылки на них. В report достаточно краткого описания и ссылок.
\end{itemize}
Если baseline result или experimentation setup не готовы, отправить короткую причину отсутствия. \\
\hline

\textbf{Артур} &
\begin{itemize}
    \item Link на PR/Issue с backend skeleton.
    \item Link на OpenAPI/Swagger, если готов.
    \item Screenshot или link на working \texttt{GET /health}, если готов.
    \item Короткая LaTeX-секция \texttt{Implemented Backend Features}: что реально успели реализовать.
\end{itemize}
Если OpenAPI, PR или \texttt{GET /health} не готовы, отправить короткую причину отсутствия для каждого отсутствующего пункта. \\
\hline

\textbf{Амир} &
\begin{itemize}
    \item Link на file/PR с initial DB schema.
    \item Короткая LaTeX-секция или ссылка \texttt{Database Schema}: основные таблицы или сущности.
    \item Если есть storage verification, дать screenshot, log или короткое описание.
\end{itemize}
Если schema или storage verification не готовы, отправить короткую причину отсутствия. \\
\hline

\textbf{Руслан} &
\begin{itemize}
    \item Link на \texttt{Dockerfile}, если готов.
    \item Link на \texttt{docker-compose.yml}, если готов.
    \item Если Docker не готов, LaTeX-секция \texttt{Run Instructions}: как запустить текущий проект.
    \item Link на PR/Issue по infrastructure/config/Git workflow, если есть.
\end{itemize}
Если Docker, compose или infrastructure PR не готовы, отправить короткую причину отсутствия. \\
\hline

\textbf{Рома} &
\begin{itemize}
    \item LaTeX-секция \texttt{Team Members and Roles}.
    \item LaTeX-секция \texttt{Communication Channels}.
    \item LaTeX-секция \texttt{Problem Statement}.
    \item LaTeX-секция \texttt{User Stories}.
    \item LaTeX-секция \texttt{Acceptance Criteria}.
    \item Wireframes/mockups as figures или links.
    \item User flow diagrams as figures или links.
    \item Screenshots/GIFs as figures или links, если готовы.
    \item LaTeX-секция \texttt{Project Links}: Kanban, PRs, Issues, API docs.
    \item LaTeX-секция \texttt{Internal Demo Notes and Next Steps}, если demo проведено.
\end{itemize}
Если screenshots, diagrams, links или demo notes не готовы, вставить в report короткую причину отсутствия. \\
\hline

\end{longtable}

\section{Template for Missing Artifacts}

Если обязательный материал отсутствует, используйте следующий шаблон:

\begin{verbatim}
\textbf{Missing artifact:} <artifact name>

\textbf{Reason:} <short reason why it is not available yet>

\textbf{Next step:} <when or how it will be completed>
\end{verbatim}

Пример для вставки в report:

\begin{verbatim}
\textbf{Missing artifact:} OpenAPI/Swagger link

\textbf{Reason:} The backend API contract is still being finalized,
so the Swagger page is not available yet.

\textbf{Next step:} The OpenAPI documentation will be added after
the issue workflow endpoints are merged.
\end{verbatim}

\section{Important Note}

В report не нужно обещать технические части, которые команда не успела реализовать. Если какой-то технический артефакт не готов, лучше честно указать причину отсутствия и next step. Это лучше, чем добавлять в report неподтверждённые features.

\end{document}
